\documentclass[12pt]{ctexart}
\usepackage{a4}
\usepackage{graphicx}
\usepackage{geometry}
\usepackage{amsmath}
\usepackage{mathtools}
\usepackage{amssymb}
\geometry{a4paper,left=2cm,right=2cm,top=2.5cm,bottom=2.5cm}

\title{几何基础作业10}
\author{无名}
\date{\today}

\begin{document}
\maketitle

\section{设$P(x_0,y_0)$是椭圆上一点，利用仿射变换求过点$P$的椭圆切线方程}

椭圆方程:$C:\frac{x^2}{a^2}+\frac{y^2}{b^2}=1$

圆方程：$C':x'^2+y'^2=1$

构造映射$f,\mathbb{R}^2 \to \mathbb{R}^2: \frac{x}{a}=x',\frac{y}{b}=y'$，显然它是一个一一映射

那么椭圆$C$经过仿射变换后就是$C'$，此时$P(x_0,y_0)$在$C'$上对应点$P'(\frac{x_0}{a},\frac{y_0}{b})$

对于单位圆$C'$，切线方程为$x_0x+y_0y=1$，
那么$P'$的切线方程就是$\frac{x_0}{a}x+\frac{y_0}{b}y=1$

进行逆仿射变换，即得到原方程$\frac{x_0x}{a^2}+\frac{y_0y}{b^2}=1$为切线方程

\section{在普通欧氏平面上$SA,SB,SC,SD$是过定点$S$的四条直线,直线$l$与这4条直线分别交于$A,B,C,D$证明$ABCD$是定值}

在$\bigtriangleup SAC$ $\bigtriangleup SBC$中，由正弦定理
\[\frac{AC}{\sin \angle ASC}=\frac{SC}{\sin \angle SAC}
\]
\[\frac{BC}{\sin \angle BSC}=\frac{SC}{\sin \angle SBC}
\]
那么
\[\frac{AC}{BC}=\frac{\sin \angle ASC\cdot \sin \angle SBC}{\sin \angle BSC \sin \angle SAC}
\]
同理
\[\frac{AD}{BD}=\frac{\sin \angle ASD\cdot \sin \angle SBD}{\sin \angle BSD \sin \angle SAD}
\]
则
\[\frac{AC}{BC}:\frac{AD}{BD}=\frac{\sin \angle ASC\cdot \sin \angle SBC}{\sin \angle BSC \sin \angle SAC}\cdot \frac{\sin \angle BSD \sin \angle SAD}{\sin \angle ASD\cdot \sin \angle SBD}
\]
化简后有
\[\frac{AC}{BC}:\frac{AD}{BD}=\frac{\sin \angle ASC}{\sin \angle BSC }\cdot \frac{\sin \angle BSD }{\sin \angle ASD}
\]

我们再取一个直线$l'$,不经过$S$且与$SA,SB,SC,SD$都相交

定义$S$到$l$的中心投影$\pi_s:l\rightarrow l'$，为$l$上一点$X$,连接$XS$与$l'$交于$X'$

那么再$\pi_s$下,$A\mapsto A',B\mapsto B',C \mapsto C',D \mapsto D'$

由于
\[\angle ASC=\angle A'SC' \hspace{2em}\angle BSD=\angle B'SD'
\]
\[\angle BSC=\angle B'SC' \hspace{2em}\angle ASD=\angle A'SD'
\]
那么
\[\frac{AC}{BC}:\frac{AD}{BD}=\frac{\sin \angle ASC}{\sin \angle BSC }\cdot \frac{\sin \angle BSD }{\sin \angle ASD}=\frac{\sin \angle A'SC'}{\sin \angle B'SC' }\cdot \frac{\sin \angle B'SD' }{\sin \angle A'SD'}=\frac{A'C'}{B'C'}:\frac{A'D'}{B'D'}
\]

故$(ABCD)$选取与$l$无关，是定值
\end{document}